\subsubsection{Bedeutung und Einsatz von LiDAR-Systemen}

LiDAR (Light Detection and Ranging) ist ein laserbasiertes System zur Erfassung dreidimensionaler Strukturen Distanzmessungen. Erstmals wurden luftgestützte LiDAR-Systeme 1964 in der Forstwirtschaft zur Messung von Baumhöhen eingesetzt. \citep{thapaForestAbovegroundBiomass2025} Bis heute wird kontinuierlich an dieser Technologie geforscht, da sie in verschiedenen Anwendungsbereichen maßgebliche Vorteile mit sich bringen. Im Gegensatz zu Kameras benötigen LiDAR-Systeme kein externes Umgebungslicht. Außerdem sind die Messungen materialunabhängig und liefern extrem schnell hochpräzise Ergebnisse. \citep{lopacApplicationLaserSystems2022}

Die Gewinnung von LiDAR-Daten erfolgt über unterschiedliche Plattformen, deren Wahl jeweils von der spezifischen Anforderungen an Reichweite und Präzision abhängt. Während stationäre Systeme wie das Terrestrial Laser Scanning (TLS) durch ihre feste Montage zwar in ihrem räumlichen Wirkungsbereich begrenzt sind, liefern sie hochpräzise Aufnahmen einzelner Objekte. Ein vergleichbar hohe Genauigkeit bieten Mobile Laser Scanning (MLS) Systeme, die für den Einsatz auf beweglichen Plattformen installiert sind und zum Beispiel im Straßenverkehr Anwendung finden, jedoch auf befahrbare Routen beschränkt bleiben. Für einen weiträumigen Überblick werden Airborne Laser Systems (ALS) in Flugzeugen oder Drohen montiert. Ihre Messgenauigkeit ist im Vergleich zu den vorherigen Verfahren etwas geringer, sie sind jedoch ideal für Anwendungen in der Fernerkundung geeignet. Den globalen Maßstab bilden Spaceborne LiDAR Systems, die zwar eine deutlich geringere Auflösung aufweisen, allerdings einen umfassenden Überblick über das gesamte Ökosystem ermöglichen. \citep{lopacApplicationLaserSystems2022}